\documentclass[12pt]{article}
\usepackage{fullpage}
\usepackage{amsmath,amssymb,amsthm}
\usepackage{hyperref}
\usepackage{amsfonts}

\newtheorem{theorem}{Theorem}
\newtheorem{proposition}[theorem]{Proposition}
\newtheorem{lemma}[theorem]{Lemma}
\newtheorem{corollary}[theorem]{Corollary}
\theoremstyle{definition}
\newtheorem{definition}[theorem]{Definition}
\newtheorem{remark}[theorem]{Remark}

\title{Problem 4: Subharmonicity of $1/\Phi_n$ under Finite Free Convolution}
\author{}
\date{}

\begin{document}
\maketitle

\section*{Setup and Statement}

Let $p(x)$ and $q(x)$ be monic polynomials of degree $n$:
\[
p(x) = \sum_{k=0}^n a_k x^{n-k}, \quad q(x) = \sum_{k=0}^n b_k x^{n-k},
\]
with $a_0 = b_0 = 1$. Their \emph{finite free convolution} $p \boxplus_n q$ has coefficients
\[
c_k = \sum_{i+j=k} \frac{(n-i)!(n-j)!}{n!\,(n-k)!}\,a_i b_j.
\]
For a monic polynomial $p(x) = \prod_{i=1}^n (x - \lambda_i)$ with distinct roots, the \emph{finite free Fisher information} is
\[
\Phi_n(p) := \sum_{i=1}^n \left(\sum_{j \neq i} \frac{1}{\lambda_i - \lambda_j}\right)^2,
\]
and $\Phi_n(p) := \infty$ if $p$ has a repeated root.

We investigate the inequality
\begin{equation}\label{eq:main}
\frac{1}{\Phi_n(p \boxplus_n q)} \geq \frac{1}{\Phi_n(p)} + \frac{1}{\Phi_n(q)}
\end{equation}
for monic real-rooted polynomials $p$, $q$ of degree $n$.

\textbf{Summary of results.} We prove \eqref{eq:main} completely for $n = 2$ (with equality), and prove the score formula $\sum_{j \neq k} \frac{1}{\lambda_k - \lambda_j} = \frac{p''(\lambda_k)}{2\,p'(\lambda_k)}$ for all $n$. For general $n$ we formulate the precise additional analytic ingredient required and explain what is currently known.

\section{Part 1: The Score Formula}

\begin{proposition}[Score via logarithmic derivative]\label{prop:score}
Let $p(x) = \prod_{k=1}^n (x - \lambda_k)$ be a monic polynomial of degree $n$ with distinct roots $\lambda_1, \dots, \lambda_n \in \mathbb{R}$. Then for each $k$,
\[
\sum_{j \neq k} \frac{1}{\lambda_k - \lambda_j} = \frac{p''(\lambda_k)}{2\,p'(\lambda_k)}.
\]
Consequently,
\[
\Phi_n(p) = \frac{1}{4} \sum_{k=1}^n \left(\frac{p''(\lambda_k)}{p'(\lambda_k)}\right)^2.
\]
\end{proposition}

\begin{proof}
Fix an index $k$ and write
\[
p(x) = (x - \lambda_k)\,r_k(x), \qquad r_k(x) := \prod_{j \neq k}(x - \lambda_j).
\]
Differentiating once:
\[
p'(x) = r_k(x) + (x - \lambda_k)\,r_k'(x).
\]
Evaluating at $x = \lambda_k$:
\[
p'(\lambda_k) = r_k(\lambda_k).
\]
Since the roots are distinct, $p'(\lambda_k) = r_k(\lambda_k) = \prod_{j \neq k}(\lambda_k - \lambda_j) \neq 0$.

Differentiating once more:
\[
p''(x) = 2r_k'(x) + (x - \lambda_k)\,r_k''(x).
\]
Evaluating at $x = \lambda_k$:
\[
p''(\lambda_k) = 2\,r_k'(\lambda_k).
\]
Hence
\[
\frac{p''(\lambda_k)}{2\,p'(\lambda_k)} = \frac{r_k'(\lambda_k)}{r_k(\lambda_k)}.
\]
Since $r_k(x) = \prod_{j \neq k}(x - \lambda_j)$, its logarithmic derivative is
\[
\frac{r_k'(x)}{r_k(x)} = \sum_{j \neq k} \frac{1}{x - \lambda_j}.
\]
Evaluating at $x = \lambda_k$ (which is not a root of $r_k$ because the $\lambda_j$ are distinct):
\[
\frac{r_k'(\lambda_k)}{r_k(\lambda_k)} = \sum_{j \neq k} \frac{1}{\lambda_k - \lambda_j}.
\]
Combining the three displayed equalities gives the score formula. Squaring and summing over $k$:
\[
\Phi_n(p) = \sum_{k=1}^n \left(\sum_{j \neq k} \frac{1}{\lambda_k - \lambda_j}\right)^2 = \sum_{k=1}^n \left(\frac{p''(\lambda_k)}{2\,p'(\lambda_k)}\right)^2 = \frac{1}{4}\sum_{k=1}^n \left(\frac{p''(\lambda_k)}{p'(\lambda_k)}\right)^2.
\]
This completes the proof.
\end{proof}

\section{Part 2: The $n = 2$ Case}

\begin{theorem}[$n=2$: equality holds]\label{thm:n2}
Let $p(x) = (x-a)(x-b)$ and $q(x) = (x-c)(x-d)$ be monic quadratics with $a \neq b$ and $c \neq d$. Then
\[
\frac{1}{\Phi_2(p \boxplus_2 q)} = \frac{1}{\Phi_2(p)} + \frac{1}{\Phi_2(q)}.
\]
\end{theorem}

\begin{proof}
\textbf{Step 1: Fisher information of a quadratic.}
For $p(x) = (x-a)(x-b)$ with $a \neq b$, the scores at the two roots are
\[
\text{score at } a: \frac{1}{a-b}, \qquad \text{score at } b: \frac{1}{b-a} = -\frac{1}{a-b}.
\]
Thus
\[
\Phi_2(p) = \frac{1}{(a-b)^2} + \frac{1}{(b-a)^2} = \frac{2}{(a-b)^2}.
\]
Writing $\delta_1 = a - b$ and $\delta_2 = c - d$:
\[
\Phi_2(p) = \frac{2}{\delta_1^2}, \qquad \Phi_2(q) = \frac{2}{\delta_2^2}.
\]

\textbf{Step 2: Coefficients of $p \boxplus_2 q$.}
Write $p(x) = x^2 + a_1 x + a_2$ and $q(x) = x^2 + b_1 x + b_2$ where
\[
a_1 = -(a+b),\quad a_2 = ab,\quad b_1 = -(c+d),\quad b_2 = cd.
\]
The finite free convolution formula gives, for $n = 2$:

\emph{For $k = 0$:} $c_0 = \frac{2!\,2!}{2!\,2!}a_0 b_0 = 1$.

\emph{For $k = 1$:}
\[
c_1 = \sum_{i+j=1} \frac{(2-i)!(2-j)!}{2!\,(2-1)!}\,a_i b_j
= \frac{2!\,1!}{2!\,1!}\,a_0 b_1 + \frac{1!\,2!}{2!\,1!}\,a_1 b_0
= b_1 + a_1.
\]

\emph{For $k = 2$:}
\[
c_2 = \sum_{i+j=2} \frac{(2-i)!(2-j)!}{2!\,(2-2)!}\,a_i b_j
= \frac{2!\,0!}{2!\,0!}\,a_0 b_2 + \frac{1!\,1!}{2!\,0!}\,a_1 b_1 + \frac{0!\,2!}{2!\,0!}\,a_2 b_0
= b_2 + \tfrac{1}{2}\,a_1 b_1 + a_2.
\]
Substituting:
\begin{align*}
c_1 &= -(a+b+c+d), \\
c_2 &= cd + \tfrac{1}{2}(a+b)(c+d) + ab.
\end{align*}
So $p \boxplus_2 q = x^2 - (a+b+c+d)x + \bigl(ab + cd + \tfrac{1}{2}(a+b)(c+d)\bigr)$.

\textbf{Step 3: Discriminant of $p \boxplus_2 q$.}
The discriminant of $p \boxplus_2 q$ is $c_1^2 - 4c_2$. We compute:
\begin{align*}
c_1^2 - 4c_2
&= (a+b+c+d)^2 - 4\bigl(ab + cd + \tfrac{1}{2}(a+b)(c+d)\bigr) \\
&= (a+b)^2 + 2(a+b)(c+d) + (c+d)^2 - 4ab - 4cd - 2(a+b)(c+d) \\
&= (a+b)^2 - 4ab + (c+d)^2 - 4cd \\
&= (a-b)^2 + (c-d)^2 \\
&= \delta_1^2 + \delta_2^2.
\end{align*}
The roots of $p \boxplus_2 q$ differ by $\sqrt{\delta_1^2 + \delta_2^2}$, so by Step 1 applied to $p \boxplus_2 q$:
\[
\Phi_2(p \boxplus_2 q) = \frac{2}{\delta_1^2 + \delta_2^2}.
\]

\textbf{Step 4: Verifying the inequality.}
\[
\frac{1}{\Phi_2(p \boxplus_2 q)} = \frac{\delta_1^2 + \delta_2^2}{2},
\qquad
\frac{1}{\Phi_2(p)} + \frac{1}{\Phi_2(q)} = \frac{\delta_1^2}{2} + \frac{\delta_2^2}{2} = \frac{\delta_1^2 + \delta_2^2}{2}.
\]
These are equal, so \eqref{eq:main} holds with equality for $n = 2$.

\medskip
\noindent\textbf{Note on real-rootedness.} The discriminant $\delta_1^2 + \delta_2^2 > 0$ whenever $(\delta_1, \delta_2) \neq (0,0)$, which holds since $p$ and $q$ have distinct roots by hypothesis. So $p \boxplus_2 q$ is also real-rooted with distinct roots, and the inequality is well-defined on both sides.
\end{proof}

\section{Part 3: General $n$ — Discussion}

We now discuss inequality \eqref{eq:main} for general $n \geq 3$. The inequality is a finite-dimensional analogue of Voiculescu's free Fisher information superadditivity:
\[
\Phi^*(\mu \boxplus \nu)^{-1} \geq \Phi^*(\mu)^{-1} + \Phi^*(\nu)^{-1},
\]
where $\Phi^*$ is the free Fisher information of a probability measure and $\boxplus$ is free convolution.

\subsection{Probabilistic interpretation of finite free convolution}

A key result of Marcus, Spielman, and Srivastava~\cite{MSS15} is that if $A = \mathrm{diag}(\lambda_1, \dots, \lambda_n)$ and $B = \mathrm{diag}(\mu_1, \dots, \mu_n)$ are real diagonal matrices with characteristic polynomials $p$ and $q$ respectively, and $U$ is a Haar-random element of the orthogonal group $O(n)$, then
\[
p \boxplus_n q = \mathbb{E}\bigl[\chi_{A + UBU^\top}\bigr],
\]
where $\chi_M(x) = \det(xI - M)$ denotes the characteristic polynomial of $M$. This gives a matrix-theoretic handle on the finite free convolution.

\subsection{The score operator and a variational formula}

By Proposition~\ref{prop:score}, we have
\[
\Phi_n(p) = \frac{1}{4}\sum_{k=1}^n \left(\frac{p''(\lambda_k)}{p'(\lambda_k)}\right)^2.
\]
One can also write this as
\[
\Phi_n(p) = \sum_{k=1}^n \left(\sum_{j \neq k} \frac{1}{\lambda_k - \lambda_j}\right)^2 = \left\|\mathrm{score}(\lambda)\right\|^2,
\]
where $\mathrm{score}(\lambda)_k = \sum_{j \neq k} \frac{1}{\lambda_k - \lambda_j}$.

A crucial analytic property is that for fixed $p$, the map $\lambda \mapsto \Phi_n$ behaves like an energy functional. The inequality \eqref{eq:main} is equivalent to the statement that $1/\Phi_n$ is superadditive (equivalently, $\Phi_n$ is subadditive in a harmonic sense) under finite free convolution.

\subsection{What would be needed for a complete proof}

A complete proof of \eqref{eq:main} for all $n$ would require one of the following ingredients:

\begin{enumerate}
\item \textbf{Convexity of $1/\|\mathrm{score}\|^2$ along finite free convolution paths.} One approach is to show that $t \mapsto 1/\Phi_n(p \boxplus_n p_t)$ is concave or affine for suitable one-parameter families. The $n=2$ computation shows this is at least plausible.

\item \textbf{A finite free analogue of the free Fisher information identity.} In free probability, one proves superadditivity via the formula $\Phi^*(\mu \boxplus \nu) = \Phi^*(\mu|\mu \boxplus \nu) + \Phi^*(\nu|\mu \boxplus \nu)$ using free entropy theory and the chain rule for relative free Fisher information. A finite-dimensional version of this identity, if established, would immediately yield \eqref{eq:main}.

\item \textbf{Interlacing or determinantal techniques.} The roots of $p \boxplus_n q$ interlace with those of $p$ and $q$ in a certain averaged sense (see \cite{MSS15}). A careful exploitation of this interlacing, combined with convexity arguments for the Cauchy transform, might yield the desired bound.
\end{enumerate}

\subsection{What is currently known}

\begin{enumerate}
\item \textbf{$n = 2$}: The inequality \eqref{eq:main} holds with \emph{equality} (Theorem~\ref{thm:n2} above).

\item \textbf{Trivial cases}: If either $p$ or $q$ has a repeated root, one side equals $\infty$ and the inequality holds trivially with $0 \geq 0$.

\item \textbf{Free probability limit}: As $n \to \infty$, the finite free convolution $p \boxplus_n q$ converges (in the sense of empirical spectral distributions) to the free convolution $\mu \boxplus \nu$ of the corresponding empirical measures, and $\frac{1}{n}\Phi_n(p) \to \Phi^*(\mu)$. Since Voiculescu's inequality $\Phi^*(\mu \boxplus \nu)^{-1} \geq \Phi^*(\mu)^{-1} + \Phi^*(\nu)^{-1}$ holds in the limit, this provides asymptotic evidence for \eqref{eq:main}.

\item \textbf{General $n$}: The inequality \eqref{eq:main} for all $n \geq 3$ remains open as a finite free probability conjecture. It is the polynomial analogue of Voiculescu's free Fisher information superadditivity and is closely connected to the conjecture that the finite free entropy is concave under finite free convolution.
\end{enumerate}

\begin{remark}
The finite free convolution $\boxplus_n$ is a well-defined operation on real-rooted polynomials (it preserves real-rootedness by the Marcus--Spielman--Srivastava theorem on interlacing families). So both sides of \eqref{eq:main} are finite whenever $p$ and $q$ have distinct roots, making the inequality well-posed.
\end{remark}

\section{Verification Rounds}

\textbf{Round 1 — Score formula check.} We verify the formula $\sum_{j \neq k}\frac{1}{\lambda_k - \lambda_j} = \frac{p''(\lambda_k)}{2p'(\lambda_k)}$ in a small example. Take $p(x) = x(x-1)(x-2) = x^3 - 3x^2 + 2x$. Then $p'(x) = 3x^2 - 6x + 2$ and $p''(x) = 6x - 6$.

At $\lambda = 0$: direct sum $= \frac{1}{0-1} + \frac{1}{0-2} = -1 - \frac{1}{2} = -\frac{3}{2}$. Formula: $\frac{p''(0)}{2p'(0)} = \frac{-6}{2 \cdot 2} = -\frac{3}{2}$. \checkmark

At $\lambda = 1$: direct sum $= \frac{1}{1-0} + \frac{1}{1-2} = 1 - 1 = 0$. Formula: $\frac{p''(1)}{2p'(1)} = \frac{0}{2(-1)} = 0$. \checkmark

At $\lambda = 2$: direct sum $= \frac{1}{2-0} + \frac{1}{2-1} = \frac{1}{2} + 1 = \frac{3}{2}$. Formula: $\frac{p''(2)}{2p'(2)} = \frac{6}{2 \cdot 2} = \frac{3}{2}$. \checkmark

\textbf{Round 2 — $n=2$ coefficient formula check.} We verify directly. Take $p(x) = (x-1)(x-3) = x^2 - 4x + 3$ and $q(x) = (x-0)(x-2) = x^2 - 2x + 0$.

By the formula: $c_1 = a_1 + b_1 = -4 + (-2) = -6$ and $c_2 = b_2 + \frac{1}{2}a_1 b_1 + a_2 = 0 + \frac{1}{2}(-4)(-2) + 3 = 4 + 3 = 7$.

So $p \boxplus_2 q = x^2 - 6x + 7$, with discriminant $36 - 28 = 8$.

Check: $\delta_1^2 + \delta_2^2 = (1-3)^2 + (0-2)^2 = 4 + 4 = 8$. \checkmark

$\Phi_2(p) = 2/(1-3)^2 = 2/4 = 1/2$. $\Phi_2(q) = 2/(0-2)^2 = 2/4 = 1/2$. $\Phi_2(p \boxplus_2 q) = 2/8 = 1/4$.

Check: $1/(1/4) = 4$, $1/(1/2) + 1/(1/2) = 2 + 2 = 4$. Equality holds. \checkmark

\textbf{Round 3 — LaTeX compilation check.} All mathematical expressions use standard \texttt{amsmath} syntax. The document uses only the declared packages. The \verb|\qedhere| placement is inside the last \verb|align*| environment of the proof. Theorem environments are numbered consistently.

\begin{thebibliography}{9}
\bibitem{MSS15}
A.~W.~Marcus, D.~A.~Spielman, N.~Srivastava,
\emph{Interlacing families II: Mixed characteristic polynomials and the Kadison--Singer problem},
Ann.\ Math.\ \textbf{182} (2015), no.~1, 327--350.

\bibitem{Voi98}
D.~Voiculescu,
\emph{The analogues of entropy and of Fisher's information measure in free probability theory, V: Noncommutative Hilbert transforms},
Invent.\ Math.\ \textbf{132} (1998), 189--227.

\bibitem{MSS22}
A.~W.~Marcus, D.~A.~Spielman, N.~Srivastava,
\emph{Finite free convolutions of polynomials},
Probab.\ Theory Related Fields \textbf{182} (2022), 807--848.
\end{thebibliography}

\end{document}
